\section{Positive completions and the moving prime boundary}\label{sec:positive}
\subsection{The unsubtracted positive prime form has zero domain}
For finite A define the positive expression
\begin{equation}
q_A^{\rm prime}[F]=
\sum_{2\le a\le A}\Lambda(a)\norm{(I-V_a)F}_\nu^2 .
\end{equation}
Its negative cross terms resemble the prime part of Q, but its
diagonal terms cannot be omitted from a positive identity.

\begin{proposition}\label{prop:prime-zero-domain}
In the full actual class sector,
\begin{equation}
\left\{F:\sum_{a\ge2}\Lambda(a)\norm{(I-V_a)F}_\nu^2<\infty\right\}
=\{0\}.
\end{equation}
\end{proposition}
\begin{proof}
Coefficient shifts and the norm equivalence \eqref{eq:gram}
give $V_aF\rightharpoonup0$ as $a\to\infty$, first for finite
characters and then for all F by uniform boundedness.
Equation \eqref{eq:weighted-adjoint} gives
\[
\norm{V_aF}_\nu^2=\int|F|^2\rho_a\dd g
\longrightarrow t_0\norm F_{\rm Haar}^2 .
\]
Consequently
$\norm{(I-V_a)F}_\nu^2\to
\norm F_\nu^2+t_0\norm F_{\rm Haar}^2$, which is strictly
positive for nonzero F. The sum already diverges along the
infinitely many primes.
\end{proof}
On raw packets the limiting positive form at finite A is
\begin{equation}
\left[\sum_{2\le a\le A}\Lambda(a)
       \left(1+\frac{\rho_a(0)}{\rho(0)}\right)\right]\norm f_2^2
-2\sum_{2\le a\le A}\frac{\Lambda(a)}{\sqrt a}
                         \Ree\ip f{U_{\log a}f}_2.
\end{equation}
The bracket diverges while, for a fixed compact f, the cross
sum stabilizes. Pole neutrality does not remove the bracket.
Formal subtraction leaves a signed expression; positivity of
the unrenormalized forms supplies no proof of positivity of
that difference.

\subsection{A coercive variational attempt also diverges}
Let h be any densely defined closed quadratic form on the class
space with $h[u]\ge\delta\norm u_\nu^2$ for $\delta>0$.
Set
\begin{equation}
q_A[u]=h[u]+q_A^{\rm prime}[u],\qquad K_A^{\rm pos}\ge\delta I .
\end{equation}
Each $q_A$ is closed on $\Dom h$. For a bounded readout
$B:\cH_{\rm cl}\to\mathcal Y$ define
\begin{equation}
\Gamma_A[y]=\inf_{Bu=y}q_A[u] .
\end{equation}
When feasible, the affine constraint is closed and strict
coercivity gives a unique minimizer. The vector
$(K_A^{\rm pos})^{1/2}u_A$ is an actual positive source.
For finite-dimensional surjective B,
\begin{equation}
u_A=(K_A^{\rm pos})^{-1}B^*
[B(K_A^{\rm pos})^{-1}B^*]^{-1}y .
\end{equation}

\begin{theorem}[Failure of a fixed bounded constrained readout]\label{thm:schur}
For every nonzero feasible y, $\Gamma_A[y]\to+\infty$.
Also $(K_A^{\rm pos})^{-1}\to0$ and
$(K_A^{\rm pos})^{-1/2}\to0$ strongly.
\end{theorem}
\begin{proof}
If a subsequence of constrained energies were bounded,
coercivity would give $u_{A_j}\rightharpoonup u$ along a
subsequence. For each fixed finite cutoff B, lower
semicontinuity of its bounded positive prime form gives
$q_B^{\rm prime}[u]\le C$. Taking $B\to\infty$ and using
Proposition~\ref{prop:prime-zero-domain} forces u=0, whereas
weak continuity of the readout forces Bu=y. This is a
contradiction.

For $u_A=(K_A^{\rm pos})^{-1}v$, coercivity bounds both
$u_A$ and $q_A[u_A]=\ip{u_A}v$. Every weak limit is zero
by the same argument without a readout. Thus
$q_A[u_A]\to0$ and $\norm{u_A}\to0$. Finally
$\norm{(K_A^{\rm pos})^{-1/2}v}^2=\ip v{u_A}\to0$.
\end{proof}
For the constraint that the Haar character coefficient $c_m=b$,
one has the explicit lower bound
\begin{equation}
q_A[u]\ge c_\rho|b|^2
\left[\delta+\sum_{\substack{p\le A\\p\ {\rm prime}}}
                       \frac{\delta\log p}{\delta+\log p}\right].
\end{equation}
Indeed keep the mass terms at m and the distinct labels pm,
and the p-th difference at pm. Minimizing
$\delta|c_{pm}|^2+\log p\,|c_{pm}-b|^2$ gives the summand.
A physical mass gap can motivate coercivity only after an
actual excited-sector identification. Even granting it here
does not repair this divergent positive construction.

\subsection{An explicit boundary layer of the signed prime operator}
The sharp arithmetic cutoff has an additional, independently
provable norm obstruction. Let $h\in\cD$ be real, supported
in $[-1,1]$, and put $f=\T h$. Define
\begin{equation}
\mathcal P_A f=\sum_{2\le a\le A}\frac{\Lambda(a)}{\sqrt a}
(U_{\log a}+U_{-\log a})f,\qquad L=\log A .
\end{equation}

\begin{theorem}[Prime-cutoff edge profiles]\label{thm:edges}
Uniformly for $s\in[-1,1]$,
\begin{align}
A^{-1/2}\mathcal P_A f(L+s)&\longrightarrow h'(s)+\tfrac12h(s),\\
A^{-1/2}\mathcal P_A f(-L+s)&\longrightarrow-h'(s)+\tfrac12h(s).
\end{align}
For the completed signed operator $L_A$ in
\eqref{eq:finite-signed},
\begin{equation}\label{eq:edge-lower}
\liminf_{A\to\infty}A^{-1}\norm{L_Af}_2^2
\ge2\norm{h'}_2^2+\tfrac12\norm h_2^2>0
\end{equation}
if h is nonzero. This uses the prime number theorem, not RH.
\end{theorem}
\begin{proof}
Let $\Psi(x)=\sum_{a\le x}\Lambda(a)$. At the positive edge
only one shift orientation contributes for large A, and its
scaled Stieltjes sum is
\[
\int_0^1 q^{-1/2} f(s-\log q)\dd\{\Psi(Aq)/A\}.
\]
The tests vanish unless $q\ge e^{-2}$. The prime number
theorem $\Psi(x)=x+o(x)$, uniformly on this fixed positive
q interval after division by A, and integration by parts
give convergence to the Lebesgue integral, uniformly in s.
Uniformity follows from bounded smooth derivatives in the
compact $(q,s)$ rectangle. The integral equals
\[
e^{s/2}\int_s^1e^{-u/2}f(u)\dd u=h'(s)+h(s)/2,
\]
because $f=-h''+h/4$ and all endpoint derivatives vanish.
The negative edge similarly gives
$e^{-s/2}\int_{-1}^se^{u/2}f(u)\dd u=-h'(s)+h(s)/2$.

Outside the support of f, the archimedean operator has kernel
$-n_\Gamma(|x-y|)$. Expanding it into exponentials shows that
its $e^{-|x|/2}$ tail vanishes by the pole moments; the
remaining tail is $O(e^{-5|x|/2})$.
It is negligible at both moving edges after scaling.
Integrate the two squared edge profiles on their disjoint
intervals. Their cross derivative terms cancel, leaving
$2\norm{h'}^2+\norm h^2/2$.
\end{proof}
Any counterterm whose output stays in a fixed compact
neighborhood of the original support, even with A-dependent
coefficients, misses these moving intervals and cannot cancel
\eqref{eq:edge-lower}. The theorem does not cover every smooth
regulator, or a nonlocal moving counterterm with an independently
justified physical definition.

This is also a lower bound in the actual state, not solely an
arithmetic calculation. For finite A let
$q_{R,A}=\UU_RK_AF_Rf$, with the signed $K_A$ of
\eqref{eq:finite-signed}. Strong recentering and its finite-A
mixed identity imply
\begin{equation}\label{eq:physical-edge}
\liminf_R\norm{q_{R,A}}_\nu\ge\norm{L_Af}_2 .
\end{equation}
To prove this, take any subsequence of bounded norms and a
weak limit q. Pair with the strongly convergent
$\UU_RF_Rg\to J_{\rm rec}g$. The limit is $\ip g{L_Af}_2$.
Density of $\cD$ and the isometry of $J_{\rm rec}$ force
the projection of q onto its range to be $J_{\rm rec}L_Af$.
If the norms are unbounded there is nothing to prove.
Combining \eqref{eq:physical-edge} and \eqref{eq:edge-lower}
gives an iterated physical norm lower bound of order $\sqrt A$.
No joint cutoff convergence is inferred.
